% !TeX root = ../main.tex
\chapter{Ereditarietà e Polimorfismo}

\section{Classi, Superclassi e Sottoclassi}

L'ereditarietà è il meccanismo che permette di creare nuove classi basandosi su classi esistenti. L'idea fondamentale è quella di riutilizzare i metodi e i campi di una classe esistente, aggiungendo o modificando funzionalità per rispondere a esigenze più specifiche.

Il criterio fondamentale per stabilire se l'ereditarietà sia appropriata è la relazione \textbf{"is-a"} (è-un). 
Ogni istanza della sottoclasse deve poter essere considerata un'istanza della superclasse. Ad esempio, poiché ogni \textit{Manager} è un \textit{Dipendente}, la classe \texttt{Manager} può ereditare da \texttt{Employee}. Richiamiamo innanzitutto la classe \texttt{Employee} che abbiamo visto nei capitoli precedenti:

\begin{lstlisting}[language=Java, caption={La superclasse Employee originale}]
import java.time.LocalDate;

public class Employee {
    private String name;
    private double salary;
    private LocalDate hireDay;

    public Employee(String name, double salary, int year, int month, int day) {
        this.name = name;
        this.salary = salary;
        this.hireDay = LocalDate.of(year, month, day);
    }

    public String getName() {
        return name;
    }

    public double getSalary() {
        return salary;
    }

    public LocalDate getHireDay() {
        return hireDay;
    }

    public void raiseSalary(double byPercent) {
        double raise = salary * byPercent / 100;
        salary += raise;
    }
}
\end{lstlisting}

In Java, l'ereditarietà si esprime tramite la parola chiave \texttt{extends}.

\begin{lstlisting}[language=Java, caption={Esempio di ereditarietà}]
public class Manager extends Employee {
    private double bonus; // Campo aggiuntivo

    public void setBonus(double bonus) {
        this.bonus = bonus;
    }
}
\end{lstlisting}

Nonostante il prefisso "super", una superclasse non è "superiore" in termini di capacità. Al contrario, la sottoclasse è solitamente più complessa:
\begin{itemize}
    \item La \textbf{Superclasse} contiene i metodi più generali.
    \item La \textbf{Sottoclasse} eredita tutto ciò che c'è nella superclasse e aggiunge campi o metodi specializzati (\textit{campi e metodi propri}).
\end{itemize}

\begin{caution}[Precisazione sull'esempio del libro]
    Nel libro di testo, l'esempio viene proposto nell'ipotesi (puramente esemplificativa) che un Employee non possa mai diventare Manager. In una situazione reale questo non è in genere realistico e la modellazione delle classi dovrebbe tenere conto di questa possibilità.
\end{caution}

\begin{observation}[Accesso di classe vs Accesso di istanza]
    In Java, il modificatore \texttt{private} garantisce la visibilità a livello di \textbf{classe} e non di singola istanza:
    \begin{itemize}
        \item Un'istanza di \texttt{Employee} può accedere ai membri privati di un'altra istanza della stessa classe.
        \item Una sottoclasse \texttt{Manager}, pur ereditando lo stato, è considerata \textbf{codice esterno} rispetto alla definizione della superclasse.
    \end{itemize}
    \textbf{Conseguenza}: La sottoclasse possiede i dati della superclasse (memoria), ma deve interagire con essi tramite l'interfaccia pubblica (metodi), preservando l'integrità dell'incapsulamento.
\end{observation}

\section{Sovrascrivere i Metodi (Overriding)}

Spesso i metodi della superclasse non sono del tutto appropriati per la sottoclasse. Nel nostro esempio, il metodo \texttt{getSalary()} di \texttt{Employee} restituisce solo lo stipendio base, mentre per un \texttt{Manager} dovrebbe restituire la somma di stipendio e bonus.
Tentare di implementare il metodo accedendo direttamente al campo della superclasse fallisce:
\begin{lstlisting}[language=Java]
public double getSalary() {
    return salary + bonus; // ERRORE: salary e privato in Employee
}
\end{lstlisting}
Anche se \texttt{Manager} è una sottoclasse, deve rispettare l'incapsulamento e usare l'interfaccia pubblica della superclasse.
Un errore altro comune è chiamare il metodo per nome sperando che Java "capisca" di voler usare quello del padre:
\begin{lstlisting}[language=Java]
public double getSalary() {
    double baseSalary = getSalary(); // ERRORE: chiama se stesso all'infinito
    return baseSalary + bonus;
}
\end{lstlisting}
Questo codice causa un \textit{crash} del programma perché il metodo richiama se stesso senza sosta (Stack Overflow).

Per istruire il compilatore a invocare il metodo della superclasse anziché quello corrente, si utilizza la parola chiave \textbf{\texttt{super}}.

\begin{lstlisting}[language=Java, caption={Overriding corretto con super}]
public double getSalary() {
    double baseSalary = super.getSalary();
    return baseSalary + bonus;
}
\end{lstlisting}

\begin{important}[\texttt{super} non è un riferimento]
A differenza di \texttt{this}, che è un riferimento all'oggetto corrente (e può essere assegnato a una variabile), \textbf{\texttt{super}} è solo una \textbf{parola chiave speciale} che dirige il compilatore verso la superclasse. Non puoi, ad esempio, scrivere \texttt{Object x = super;}.
\end{important}

\subsection{L'importanza di \texttt{@Override}}
Le annotazioni in Java sono metadati che forniscono informazioni aggiuntive al compilatore; saranno trattate più approfonditamente più avanti, ma è importante introdurre subito l'annotazione \texttt{@Override} in questo contesto.
Sebbene non sia strettamente necessaria per il funzionamento del codice, è buona norma precedere ogni metodo sovrascritto con l'annotazione \textbf{\texttt{@Override}}.

\begin{lstlisting}[language=Java, caption={Uso di @Override}]
@Override
public double getSalary() {
    return super.getSalary() + bonus;
}
\end{lstlisting}

[Motivazioni e Vantaggi]
L'uso di questa annotazione serve a istruire il compilatore sulle nostre intenzioni. Le motivazioni principali sono:
\begin{enumerate}
    \item \textbf{Controllo degli errori di battitura}: Se scrivessimo per errore \texttt{getsalary()} (con la 's' minuscola), senza annotazione Java penserebbe semplicemente che stiamo creando un \textit{nuovo} metodo. Con \texttt{@Override}, il compilatore segnalerebbe un errore perché non esiste alcun metodo \texttt{getsalary} nella superclasse da sovrascrivere.
    \item \textbf{Verifica della firma}: Se cambiamo i parametri del metodo (es. aggiungiamo un argomento), non stiamo più facendo \textit{overriding} ma \textit{overloading}. L'annotazione ci avvisa immediatamente se stiamo rompendo il legame con il metodo del padre.
    \item \textbf{Documentazione}: Rende immediatamente chiaro a chi legge il codice che quel metodo non è un'aggiunta originale della sottoclasse, ma una modifica di un comportamento ereditato.
\end{enumerate}



\begin{tip}[Fail-Fast]
Usare \texttt{@Override} è un esempio del principio \textit{"fail-fast"}: è meglio che il compilatore ci urli contro subito per un errore di distrazione, piuttosto che scoprire a runtime che il nostro \texttt{Manager} sta prendendo lo stipendio base perché il metodo di calcolo del bonus ha un nome leggermente sbagliato.
\end{tip}

\section{I Costruttori delle Sottoclassi}

A differenza dei metodi, i costruttori non vengono "ereditati" nel senso stretto del termine: ogni sottoclasse deve definire i propri costruttori, ma ha l'obbligo di invocare un costruttore della superclasse per inizializzare i campi privati ereditati. Il meccanismo di \textit{constructor chaining} (catena dei costruttori) impone infatti dei vincoli precisi a seconda della struttura della superclasse. Si possono verificare due scenari principali:

\begin{enumerate}
    \item \textbf{Superclasse con costruttore no-args (esplicito o implicito)}
    Se nella superclasse esiste un costruttore senza argomenti, la sottoclasse non è obbligata a dichiarare un costruttore proprio. In questo caso, viene utilizzato il costruttore di default della sottoclasse che invoca implicitamente \texttt{super()}.

    \item \textbf{Superclasse senza costruttore no-args}
    Se nella superclasse è presente almeno un costruttore con parametri e \textbf{non} è stato definito un costruttore senza argomenti, la sottoclasse ha un obbligo formale:
    bisogna dichiarare esplicitamente un costruttore nella sottoclasse e la prima istruzione di tale costruttore deve essere una chiamata esplicita a uno dei costruttori della superclasse tramite la sintassi \texttt{super(...parametri...)}.
\end{enumerate}

\begin{caution}[Errore di compilazione comune]
    Se la superclasse richiede parametri e la sottoclasse non li fornisce esplicitamente (o tenta di farlo dopo la prima riga), il compilatore segnalerà l'errore: \textit{"Implicit super constructor is undefined. Must explicitly invoke another constructor"}.
\end{caution}

\begin{observation}[Invocazione Implicita e Campi Propri]
    Se nella superclasse non è definito alcun costruttore esplicito (ovvero è presente il costruttore \textit{no-args} di default), la sottoclasse gode di maggiore flessibilità: è possibile definire un costruttore che inizializza esclusivamente i \textbf{campi propri} della sottoclasse. In questo scenario, la chiamata a \texttt{super()} può essere omessa (e di solito è considerata una buona pratica, per ridurre la verbosità), poiché il compilatore la inserirà automaticamente come prima riga invisibile.
\end{observation}


\begin{observation}[La regola della prima posizione]
    Entrambe le istruzioni, \texttt{this()} e \texttt{super()}, hanno il vincolo categorico di dover essere la \textbf{prima istruzione} del costruttore. Di conseguenza, \textbf{non possono mai comparire nello stesso costruttore}.
\end{observation}


\begin{observation}
Cosa succede se un costruttore nella sottoclasse usa \texttt{this(...)}? La responsabilità di invocare il padre viene semplicemente "passata" lungo la catena:
\begin{enumerate}
    \item Il Costruttore A chiama \texttt{this(...)} (Costruttore B).
    \item Il Costruttore B deve, a sua volta, chiamare \texttt{super(...)} (o un altro costruttore della stessa classe che alla fine lo faccia).
\end{enumerate}

Indipendentemente da quanti \texttt{this()} vengano usati, l'ultimo anello della catena di costruttori di una classe \textbf{deve} sempre invocare (esplicitamente o implicitamente) un costruttore della superclasse. La "nascita" del padre non può essere evitata, può solo essere delegata a un altro costruttore "fratello".
\end{observation}

\section{Polimorfismo}

Il polimorfismo è la capacità di un oggetto di essere trattato come un'istanza della sua superclasse. Questo concetto si basa sul \textbf{Principio di Sostituzione}.

Il principio afferma che è sempre possibile utilizzare un oggetto di una sottoclasse in qualsiasi contesto in cui sia attesa una superclasse. Questo perché, per la relazione \textit{is-a}, una sottoclasse possiede tutte le caratteristiche della superclasse.

\begin{lstlisting}[language=Java, caption={Esempio di sostituzione}]

// Sintassi generale:
// SuperClasse oggetto = new SottoClasse(...);

Employee e = new Manager("Rossi", 50000, 2023, 10, 1);
\end{lstlisting}

Nella riga di codice polimorfica di sopra, dobbiamo distinguere due aspetti:
\begin{itemize}
    \item \textbf{Tipo Dichiarato}: È il tipo della variabile (\texttt{Employee}). Determina quali metodi il compilatore ci permetterà di chiamare: \textit{non è possibile invocare metodi specifici} di \texttt{Manager} su una variabile dichiarata come \texttt{Employee}; se provassimo a chiamare \texttt{e.setBonus(10000)}, otterremmo un errore di compilazione.
    \item \textbf{Tipo Effettivo}: È il tipo dell'oggetto creato in memoria (\texttt{Manager}). Determina quale versione del metodo verrà eseguita a runtime.
\end{itemize}



\begin{important}[L'unidirezionalità della sostituzione]
    La sostituzione funziona solo verso l'alto nella gerarchia:
    \begin{itemize}
        \item \textbf{OK}: Assegnare un \texttt{Manager} a una variabile \texttt{Employee}.
        \item \textbf{ERRORE}: Assegnare un \texttt{Employee} a una variabile \texttt{Manager}.
    \end{itemize}
\end{important}

\begin{observation}[Array Polimorfici]
    Il principio di sostituzione si applica anche agli array. È possibile creare un array di una superclasse e riempirlo con oggetti di sottoclassi diverse:
    \begin{lstlisting}[language=Java]
    Employee[] staff = new Employee[3];
    staff[0] = new Manager("Boss", 80000, 2020, 1, 1);
    staff[1] = new Employee("Worker", 40000, 2021, 5, 1);
    \end{lstlisting}
    Questo permette di gestire gruppi eterogenei di oggetti con un unico ciclo di elaborazione.
\end{observation}

\subsection{Il Legame Dinamico (Dynamic Binding)}

Il cuore pulsante del polimorfismo è il \textbf{dynamic binding}. È il meccanismo per cui la JVM decide, solo al momento dell'esecuzione, quale versione di un metodo invocare.

Consideriamo un array di tipo \texttt{Employee} che contiene sia dipendenti "standard" che manager:

\begin{lstlisting}[language=Java, caption={Esempio di Dynamic Binding}]
Employee[] staff = new Employee[3];

// Sostituzione: Manager "travestito" da Employee
staff[0] = new Manager("Carl Cracker", 80000, 1987, 12, 15);
staff[1] = new Employee("Harry Hacker", 50000, 1989, 10, 1);
staff[2] = new Employee("Tony Tester", 40000, 1990, 3, 15);

// Ciclo polimorfico
for (Employee e : staff) {
    System.out.println(e.getName() + " " + e.getSalary());
}
\end{lstlisting}

Quando il programma arriva alla riga \texttt{e.getSalary()}, avviene la magia del dynamic binding:
\begin{enumerate}
    \item Il \textbf{Compilatore} controlla il tipo dichiarato (\texttt{Employee}) e verifica che esista un metodo \texttt{getSalary()}. Se esiste, dà il via libera.
    \item La \textbf{JVM (a runtime)} guarda l'oggetto reale a cui punta \texttt{e}:
    \begin{itemize}
        \item Se \texttt{e} punta a un \texttt{Employee}, viene eseguito \texttt{Employee.getSalary()} (quello base).
        \item Se \texttt{e} punta a un \texttt{Manager}, viene eseguito \texttt{Manager.getSalary()} (il metodo sovrascritto, che aggiunge il bonus).

    \end{itemize}
\end{enumerate}



\begin{important}[Flessibilità ed Estensibilità]
    La bellezza del dynamic binding è che se domani aggiungessimo una nuova sottoclasse \texttt{Executive} con il suo metodo \texttt{getSalary()} specifico, il ciclo \texttt{for} qui sopra \textbf{non cambierebbe di una virgola}. Il codice diventa "aperto" a nuove forme senza dover essere riscritto.
\end{important}

\begin{observation}
    Il polimorfismo è la capacità di una variabile di riferirsi a tipi diversi. Il dynamic binding è la capacità della JVM di scegliere il metodo corretto per quel tipo specifico. Sono due facce della stessa medaglia.
\end{observation}

\subsection{Sintesi Mnemonica: Il Senso di Marcia}

Per evitare confusione nel gestire le gerarchie, possiamo definire una regola basata sulla "direzione" in cui Java percorre l'albero delle classi:

\begin{description}
    \item[\textbf{I Costruttori partono dall'alto (Top-Down)}:] 
    Quando crei un oggetto, Java garantisce che le fondamenta siano solide prima di costruire i piani superiori. L'esecuzione risale fino a \texttt{Object} e poi scende: \textit{Padre $\rightarrow$ Figlio}. Non puoi esistere come "Manager" se prima non sei stato inizializzato come "Employee".
    
    \item[\textbf{I Metodi partono dal basso (Bottom-Up)}:] 
    Quando invochi un metodo su un oggetto polimorfico, Java cerca l'implementazione più specifica possibile. Parte dalla classe effettiva dell'oggetto e risale solo se non trova una sovrascrittura: \textit{Figlio $\rightarrow$ Padre}. Si predilige sempre la "personalità" della sottoclasse rispetto alla generalità della superclasse.
\end{description}



\begin{tip}[Perché questo dualismo?]
    Questa differenza esiste per un motivo preciso:
    \begin{itemize}
        \item I \textbf{costruttori} servono alla \textbf{sicurezza}: garantiscono che ogni parte dell'oggetto sia configurata correttamente.
        \item I \textbf{metodi} servono alla \textbf{flessibilità}: permettono a una sottoclasse di cambiare comportamento senza che chi usa l'oggetto debba accorgersene.
    \end{itemize}
\end{tip}
